\subsection{问题二：住院患者48小时排程模型}

\subsubsection{患者级离散时间模型}

对住院事件$e$，项目$j\in\J_e$需要$p_{ej}$个5分钟时隙，可用设备集合为$\R_{ej}$。若项目在设备$r$、时隙$t$开始，令$x_{ejrt}=1$。每个项目至多开始一次：
\begin{equation}
  \sum_{r\in\R_{ej}}\sum_{t\in\T_{ej}}x_{ejrt}\leq1,
  \qquad e\in\Ecal,\;j\in\J_e,
  \label{eq:once}
\end{equation}
其中$\T_{ej}$已剔除释放前、截止后、跨午休和跨日的开始时隙。事件是否全部完成由
\begin{equation}
  n_e y_e\leq\sum_{j\in\J_e}\sum_{r,t}x_{ejrt},
  \qquad y_e\in\{0,1\}
  \label{eq:complete}
\end{equation}
连接；因此多项目患者缺少任一项目时$y_e=0$。

设备互斥约束为
\begin{equation}
  \sum_{e,j}\sum_{t':\,t'\leq t<t'+p_{ej}}x_{ejrt'}\leq1,
  \qquad r\in\R,\;t\in\T.
  \label{eq:room}
\end{equation}
医生是所有机房共享的并发资源：
\begin{equation}
  \sum_{e,j,r}\sum_{t':\,t'\leq t<t'+p_{ej}}x_{ejrt'}\leq c_{w(t),s(t)},
  \qquad t\in\T,
  \label{eq:doctor}
\end{equation}
其中$w(t)$为星期，$s(t)$为日内槽。式\eqref{eq:doctor}防止把同一批医生分别复制给22台设备。

\subsubsection{患者内部顺序与准备约束}

同一患者的项目必须串行。若项目$j$在$j'$之前，记$o_{ejj'}=1$，则
\begin{align}
  S_{ej'}&\geq S_{ej}+p_{ej}+\tau v_{ejj'}-M(1-o_{ejj'}),\\
  S_{ej}&\geq S_{ej'}+p_{ej'}+\tau v_{ej'j}-Mo_{ejj'},
  \label{eq:precedence}
\end{align}
其中$v_{ejj'}=1$表示相邻项目不在同一设备，$\tau=2$个时隙。对$n_e>2$的事件，精确整数模型还需引入序位变量消除子环。大规模求解中直接构造一个完整项目序列，天然满足式\eqref{eq:precedence}。

空腹项目的完成时刻满足
\begin{equation}
  S_{ej}+p_{ej}\leq 10{:}00,
  \qquad f_{ej}=1.
\end{equation}
憋尿事件的有效释放时刻为$a'_e=a_e+\delta_e$，主情景$\delta_e=60$分钟；床旁项目令$\R_{ej}=\{7\}$。其他项目不得因为患者包含床旁标记而自动绕过项目兼容，每个项目仍按自身能力集合安排。

\subsubsection{目标函数}

问题二的首要目标是最大化48小时内完成全部项目的患者数：
\begin{equation}
  \max Z_1=\sum_{e\in\Ecal_I}y_e.
  \label{eq:p2-primary}
\end{equation}
在$Z_1$固定后，依次最小化首检等待、跨室次数和设备负荷不均衡：
\begin{align}
  \min Z_2&=\sum_{e:y_e=1}(S_e-a'_e),\\
  \min Z_3&=\sum_{e,j,j'}v_{ejj'},\\
  \min Z_4&=\sum_{r\in\R}\left|L_r-\frac{1}{|\R|}\sum_{r'}L_{r'}\right|.
\end{align}
采用词典序而非$Z_1-\alpha Z_2-\beta Z_3$，是为了避免任意权重让若干分钟等待抵消一个患者的48小时成功。带患者优先级和时限保证的动态诊断资源调度也常采用服务水平优先的结构\cite{patrick2008,bauerhenne2026}。

\subsubsection{大规模构造算法}

全年事件和项目数使全时域MILP规模过大。本文保留上述整数模型作为结构定义，再采用与其共享可行域的确定性构造算法。

对每个事件，先计算项目的可用设备数$|\R_{ej}|$和工作量$p_{ej}$，形成四个候选项目序列：空腹优先且稀缺优先、空腹优先且长任务优先、纯稀缺优先、原始项目序。若所有项目存在共同设备，先在累计负荷最低的三个共同设备中依次寻找完整方案；共同设备快捷搜索不可行时，逐个候选序列在\textbf{全部}兼容设备中寻找最早可行时隙，并在换室处插入10分钟转运。因此“三个”只是全年构造算法的同室优先搜索宽度，不是现实中的设备数量限制，也不从可行集合中删除其他设备。只有事件全部项目都找到时才提交占用；任一项目失败则释放本次试排的全部槽，不产生“半个患者已排”的状态。

设备选择不仅看当前负荷，还定义背景需求稀缺权重
\begin{equation}
  \kappa_r=\sum_{e,j:\,r\in\R_{ej}}\frac{p_{ej}}{|\R_{ej}|}.
\end{equation}
式中求和只使用2019年3月1日至2024年3月31日训练期形成的项目目录与时长，不读取留出期项目构成。在同一最早时刻有多台设备可行时，先保留患者上一个房间，再选累计负荷低、$\kappa_r$较小的设备；这使灵活任务尽量不占用稀缺设备。

\subsubsection{策略定义}

比较三种策略：

\begin{enumerate}
  \item \textbf{共享FCFS：}门诊/体检按计划日和释放顺序排入，住院按释放时刻排序；设备先取最早可行时刻，同刻再按编号选择。
  \item \textbf{松弛度保护：}背景部分与共享FCFS相同；住院按$\sigma_e=d_e-p_e$排序，其中$p_e=\sum_jp_{ej}$，优先处理剩余余量小的多项目患者。
  \item \textbf{联合稀缺度：}背景事件先按计划日，再按空腹、最小可用设备数和工作量排序；住院采用松弛度排序，具体设备按负荷与稀缺度选择。
\end{enumerate}

三种策略使用完全相同的事件、特殊病例、时长、设备集合和医生容量，唯一差异是排序与同刻设备择优。因此策略差异可以归因于排程规则，而不是输入变化。

\subsubsection{问题二结果与推荐方案}

在问题三的门诊/体检共享压力下，三种策略结果见表\ref{tab:policy}和图\ref{fig:policy}. 单纯松弛度保护的住院48小时率为63.75\%，低于共享FCFS的63.95\%，说明提前放置长且稀缺的多项目患者会增加局部碎片。联合稀缺度把背景预约日率由59.70\%提高到61.61\%，但住院率降至63.31\%；它与共享FCFS互不支配，而不是双目标占优。

\begin{table}[H]
  \centering
  \caption{主情景三种患者级策略的留出期结果}
  \label{tab:policy}
  \begin{tabular}{lrrrrr}
\toprule
策略 & 住院48h完成数 & 住院率 & 背景完成数 & 背景率 & 等待P90/h \\
\midrule
\textbf{共享 FCFS} & 45,260 & 63.95\% & 57,341 & 59.70\% & 40.42 \\
松弛度保护 & 45,116 & 63.75\% & 57,341 & 59.70\% & 40.25 \\
联合稀缺度 & 44,805 & 63.31\% & 59,175 & 61.61\% & 41.50 \\
\bottomrule
\end{tabular}

\end{table}

\begin{figure}[H]
  \centering
  \includegraphics[width=0.58\textwidth]{fig04_policy_tradeoff.pdf}
  \caption{共享资源下三种排程策略比较}
  \label{fig:policy}
\end{figure}

\newpage
共享FCFS成功排入45\,260个留出期住院事件，未分配25\,511个。成功者首检等待中位数4.50小时、P90为40.42小时、P95为45.58小时；多项目患者跨室率为19.96\%。它比松弛度保护多完成144个住院事件；相对联合稀缺度则多完成455个住院事件，但少完成1\,834个门诊/体检计划日事件。由于赛题把住院48小时率设为唯一考核指标，且共享FCFS的背景率就是本文定义的不劣基线，最终推荐资源可行FCFS；若医院把背景服务下限提高到61.61\%附近，则应切换到联合稀缺度帕累托方案，而不能声称一套策略同时改善所有指标。

表\ref{tab:recommendation}给出推荐清单的开头部分。完整清单逐事件给出患者序位、项目序位、项目族、开始与结束时刻、设备、准备标记、转运时间和兼容证据层，可直接转换为预约系统导入表。

\begin{table}[H]
  \centering
  \caption{留出期住院患者推荐排程样例}
  \label{tab:recommendation}
  \begin{tabularx}{\textwidth}{r l r X l l}
\toprule
序位 & 事件 & 项目序位 & 项目族 & 时段 & 设备 \\
\midrule
1 & IP2024-000259 & 1 & 腹部/胃肠 & 04-01 08:30--08:40 & 门16(65) \\
2 & IP2024-000261 & 1 & 浅表器官 & 04-01 08:30--08:40 & 妇门2(75) \\
3 & IP2024-000265 & 1 & 浅表器官 & 04-01 08:35--08:45 & 门9(47) \\
4 & IP2024-000268 & 1 & 血管 & 04-01 08:50--09:00 & 病2(71) \\
4 & IP2024-000268 & 2 & 血管 & 04-01 09:40--09:50 & 病2(71) \\
5 & IP2024-000267 & 1 & 浅表器官 & 04-01 09:00--09:10 & 病3(72) \\
6 & IP2024-000271 & 1 & 血管 & 04-01 09:00--09:10 & 门3(42) \\
6 & IP2024-000271 & 2 & 介入/定位 & 04-01 09:45--09:55 & 门3(42) \\
\bottomrule
\end{tabularx}

\end{table}
